\documentclass{article}
\usepackage{graphicx}
\begin{document}
\begin{abstract}
We propose SparseLong, a sparse attention method. We show that it improves F1 by 3.2\% over Longformer on NQ-L.
\end{abstract}
\section{Introduction}
Long-document models are costly \cite{vaswani2017attention}. Transformers have become the dominant architecture, as shown by Vaswani et al. \cite{vaswani2017attention,ghost2099}.
% 99.9 hidden in a comment should be ignored
\section{Experiments}
Table~\ref{tab:main} reports results. SparseLong achieves 87.4\% exact match on NQ-L, outperforming BigBird at 85.1\%.
Our method reaches 41.7 ROUGE-L on NarrativeQA (Figure~\ref{fig:speed}). It also achieves a 99.2\% F1, which is state-of-the-art.
We use 12 layers and 768 hidden units.
\begin{table}
\caption{Main results.}
\label{tab:main}
\begin{tabular}{lcc}
Method & EM & F1 \\
SparseLong & 87.4 & 91.0 \\
BigBird & 85.1 & 88.2 \\
\end{tabular}
\end{table}
\begin{figure}
\includegraphics[width=0.9\linewidth]{figures/fig_speed.png}
\caption{Speed.}
\label{fig:speed}
\end{figure}
\section{Conclusion}
We demonstrate that sparse attention is effective.
\end{document}
